% Global CSVs have quoted headers, so columns are referenced by index: 0 Time, 1 kinetic_energy_int, 2 psie_active_int, 3 total_energy.
% The first row of run1_global.csv is an all-zero placeholder at the restart time, so it is dropped.
% The zoomed panel puts kinetic (solid, left axis) and elastic (dashed, right axis) energy on separate y axes;
% it has no legend of its own and shares the line styles of the full-history panel.
\pgfplotsset{
    ssbskipfirst/.style={x filter/.expression={\coordindex==0 ? nan : x}},
    ssbzoom/.style={
            width=0.8\linewidth,
            height=0.65\linewidth,
            xmin=18.5, xmax=20,
            restrict x to domain=18.5:20,
        },
}
\begin{figure}[htbp]
    \centering
    \begin{subfigure}[t]{0.48\textwidth}
        \centering
        \begin{tikzpicture}
            \begin{axis}[
                    width=0.8\linewidth,
                    height=0.65\linewidth,
                    xmin=0, xmax=20,
                    xlabel={Time},
                    ylabel={Energy},
                    legend pos=south east,
                    legend cell align=left,
                    legend style={font=\scriptsize},
                    grid=major,
                ]
                \addplot[blue, line width=1.6pt] table[col sep=comma, x index=0, y index=1] {results/data/simpleshearband/run0_global.csv};
                \addlegendentry{Kinetic, run 0}
                \addplot[red, line width=0.8pt, ssbskipfirst] table[col sep=comma, x index=0, y index=1] {results/data/simpleshearband/run1_global.csv};
                \addlegendentry{Kinetic, run 1}
                \addplot[blue, line width=1.6pt, dashed] table[col sep=comma, x index=0, y index=2] {results/data/simpleshearband/run0_global.csv};
                \addlegendentry{Elastic, run 0}
                \addplot[red, line width=0.8pt, dashed, ssbskipfirst] table[col sep=comma, x index=0, y index=2] {results/data/simpleshearband/run1_global.csv};
                \addlegendentry{Elastic, run 1}
            \end{axis}
        \end{tikzpicture}
        \caption{Full history}\label{fig:ssb_energy_full}
    \end{subfigure}
    \hfill
    \begin{subfigure}[t]{0.48\textwidth}
        \centering
        \begin{tikzpicture}
            \begin{axis}[
                    ssbzoom,
                    axis y line*=left,
                    xlabel={Time},
                    ylabel={Kinetic energy},
                    grid=major,
                ]
                \addplot[blue, line width=1.6pt] table[col sep=comma, x index=0, y index=1] {results/data/simpleshearband/run0_global.csv};
                \addplot[red, line width=0.8pt, ssbskipfirst] table[col sep=comma, x index=0, y index=1] {results/data/simpleshearband/run1_global.csv};
            \end{axis}
            \begin{axis}[
                    ssbzoom,
                    axis y line*=right,
                    axis x line=none,
                    ylabel={Elastic energy},
                    yticklabel style={/pgf/number format/fixed, /pgf/number format/precision=3},
                ]
                \addplot[blue, line width=1.6pt, dashed] table[col sep=comma, x index=0, y index=2] {results/data/simpleshearband/run0_global.csv};
                \addplot[red, line width=0.8pt, dashed, ssbskipfirst] table[col sep=comma, x index=0, y index=2] {results/data/simpleshearband/run1_global.csv};
            \end{axis}
        \end{tikzpicture}
        \caption{Zoomed near restart; kinetic (solid, left axis), elastic (dashed, right axis)}\label{fig:ssb_energy_zoom}
    \end{subfigure}
    \caption{Kinetic and elastic energy vs.\ time for run 0 and run 1.}\label{fig:ssb_energy}
\end{figure}
